% Auto-generated from book/A_terminology_notation.md — do not edit by hand
% Regenerated by scripts/md_to_latex.py

\chapter{Terminology and Notation}
\label{ch:app_a_terminology}

Nonstandard terms and the minimal notation sheet used throughout the book. (Expanded from the Preface and \texttt{HOW\_TO\_USE.md}.)

\bigskip
\noindent\rule{\textwidth}{0.4pt}
\bigskip

\section{A.1 Claim labels}
\label{ch:app_a_terminology:a-1-claim-labels}

\begin{center}
\small
\begin{tabular}{@{}lp{0.55\textwidth}@{}}
\toprule
Label & Meaning \\
\midrule
\textbf{Theorem} & Proved here or classical \\
\textbf{Model} & Consistent construction; not claimed as nature's unique law \\
\textbf{Hypothesis} & Observational claim needing validation or falsification \\
\textbf{Software fact} & True of the current code or portal implementation \\
\bottomrule
\end{tabular}
\end{center}


\bigskip
\noindent\rule{\textwidth}{0.4pt}
\bigskip

\section{A.2 Nonstandard terminology}
\label{ch:app_a_terminology:a-2-nonstandard-terminology}

\begin{center}
\small
\begin{tabular}{@{}lp{0.55\textwidth}@{}}
\toprule
Term & Working meaning \\
\midrule
\textbf{Gauged Hopf lattice} & Discrete point set in \(S^3\) (or dynamics on site quaternions) with Hopf projection, adjacency, and left/right gauge actions \\
\textbf{Flux flywheel} & Closed, topologically protected circulating flux configuration on the lattice \\
\textbf{Flux topograph} & Value landscape of a flux functional on lattice sites/edges, with separators (lift of Conway topographs) \\
\textbf{Magic Island} & Region of enhanced stability / periodicity in parameter or \(Z\) space (Model) \\
\textbf{Porous vacuum} & Informal name for incomplete / gauge-flexible ambient lattice substrate \\
\textbf{Separator structure} & Locus where a flux functional changes sign or crosses a threshold \\
\textbf{Class-group analogue} & Equivalence classes of reduced topographs/flywheels with a candidate composition law (Model; OP6) \\
\textbf{Class-number analogue} & Count of inequivalent reduced representatives (\texttt{class\_number\_analogue}) \\
**\(Z\mapsto\) map** & Portal map from atomic number to flywheel metrics (\texttt{map\_z\_to\_flywheel}) \\
\textbf{Table T4} & Pre-registered validation checklist for Hypotheses (Ch. 10, Appendix D) \\
\textbf{QVPIC} & Quaternion vortex persistent identity (related Kingdom Come / QVPIC work; optional cross-reference) \\
**\(W_g\)** & Working topological clock \(350/\pi \approx 111.408\) (Hypothesis layer) \\
\bottomrule
\end{tabular}
\end{center}


\bigskip
\noindent\rule{\textwidth}{0.4pt}
\bigskip

\section{A.3 Notation sheet}
\label{ch:app_a_terminology:a-3-notation-sheet}

\begin{center}
\small
\begin{tabular}{@{}lp{0.55\textwidth}@{}}
\toprule
Symbol & Meaning \\
\midrule
\(\mathbb{H}\) & Hamilton quaternions \\
\(S^3\) & Unit quaternions \\
\(h:S^3\to S^2\) & Hopf map \\
\(N(q)=q\overline{q}\) & Multiplicative (squared) norm \\
\(\mathcal{H}\) & Hurwitz order \\
\(L=\mathbb{Z}[i,j,k]\) & Lipschitz order \\
\(\Lambda_0\) & \(\mathcal{H}\cap S^3\) (24 units) \\
\(\Lambda_{\mathrm{ang}}\) & Angle-sampled lattice \\
\(\Lambda_{\mathrm{dyn}}\) & Two-gyro dynamical sites \\
\(E_\parallel, E_\perp\) & Along-fiber / inter-fiber edges \\
\(\bigl(\frac{a,b}{\mathbb{Q}}\bigr)\) & Quaternion algebra over \(\mathbb{Q}\) \\
OP1--OP6 & Open problems (Appendix B) \\
\bottomrule
\end{tabular}
\end{center}


\bigskip
\noindent\rule{\textwidth}{0.4pt}
\bigskip

\section{A.4 Software path cheat sheet}
\label{ch:app_a_terminology:a-4-software-path-cheat-sheet}

\begin{center}
\small
\begin{tabular}{@{}lp{0.55\textwidth}@{}}
\toprule
Path & Role \\
\midrule
\texttt{lib/hopf\_lattice.py} & Lattice, adjacency (OP1), gauge sequences \\
\texttt{lib/flux\_topograph.py} & Topographs, classification (OP2--OP3) \\
\texttt{lib/composition.py} & Composition, class-group analogue (OP6) \\
\texttt{lib/quaternion\_algebra.py} & Algebras, Hurwitz, toy class number \\
\texttt{lib/validation.py} & Table T4, OP5 diagnostics \\
\texttt{kingdom.core.*} & Live portal modules \\
\texttt{kingdom.simulations.lattice} & \texttt{TwoGyroLattice} \\
\bottomrule
\end{tabular}
\end{center}


\bigskip
\noindent\rule{\textwidth}{0.4pt}
\bigskip


